\documentclass[11pt]{article}

    \usepackage[breakable]{tcolorbox}
    \usepackage{parskip} % Stop auto-indenting (to mimic markdown behaviour)
    

    % Basic figure setup, for now with no caption control since it's done
    % automatically by Pandoc (which extracts ![](path) syntax from Markdown).
    \usepackage{graphicx}
    % Maintain compatibility with old templates. Remove in nbconvert 6.0
    \let\Oldincludegraphics\includegraphics
    % Ensure that by default, figures have no caption (until we provide a
    % proper Figure object with a Caption API and a way to capture that
    % in the conversion process - todo).
    \usepackage{caption}
    \DeclareCaptionFormat{nocaption}{}
    \captionsetup{format=nocaption,aboveskip=0pt,belowskip=0pt}

    \usepackage{float}
    \floatplacement{figure}{H} % forces figures to be placed at the correct location
    \usepackage{xcolor} % Allow colors to be defined
    \usepackage{enumerate} % Needed for markdown enumerations to work
    \usepackage{geometry} % Used to adjust the document margins
    \usepackage{amsmath} % Equations
    \usepackage{amssymb} % Equations
    \usepackage{textcomp} % defines textquotesingle
    % Hack from http://tex.stackexchange.com/a/47451/13684:
    \AtBeginDocument{%
        \def\PYZsq{\textquotesingle}% Upright quotes in Pygmentized code
    }
    \usepackage{upquote} % Upright quotes for verbatim code
    \usepackage{eurosym} % defines \euro

    \usepackage{iftex}
    \ifPDFTeX
        \usepackage[T1]{fontenc}
        \IfFileExists{alphabeta.sty}{
              \usepackage{alphabeta}
          }{
              \usepackage[mathletters]{ucs}
              \usepackage[utf8x]{inputenc}
          }
    \else
        \usepackage{fontspec}
        \usepackage{unicode-math}
    \fi

    \usepackage{fancyvrb} % verbatim replacement that allows latex
    \usepackage[Export]{adjustbox} % Used to constrain images to a maximum size
    \adjustboxset{max size={0.9\linewidth}{0.9\paperheight}}

    % The hyperref package gives us a pdf with properly built
    % internal navigation ('pdf bookmarks' for the table of contents,
    % internal cross-reference links, web links for URLs, etc.)
    \usepackage{hyperref}
    % The default LaTeX title has an obnoxious amount of whitespace. By default,
    % titling removes some of it. It also provides customization options.
    \usepackage{titling}
    \usepackage{longtable} % longtable support required by pandoc >1.10
    \usepackage{booktabs}  % table support for pandoc > 1.12.2
    \usepackage{array}     % table support for pandoc >= 2.11.3
    \usepackage{calc}      % table minipage width calculation for pandoc >= 2.11.1
    \usepackage[inline]{enumitem} % IRkernel/repr support (it uses the enumerate* environment)
    \usepackage[normalem]{ulem} % ulem is needed to support strikethroughs (\sout)
                                % normalem makes italics be italics, not underlines
    \usepackage{mathrsfs}
    

    
    % Colors for the hyperref package
    \definecolor{urlcolor}{rgb}{0,.145,.698}
    \definecolor{linkcolor}{rgb}{.71,0.21,0.01}
    \definecolor{citecolor}{rgb}{.12,.54,.11}

    % ANSI colors
    \definecolor{ansi-black}{HTML}{3E424D}
    \definecolor{ansi-black-intense}{HTML}{282C36}
    \definecolor{ansi-red}{HTML}{E75C58}
    \definecolor{ansi-red-intense}{HTML}{B22B31}
    \definecolor{ansi-green}{HTML}{00A250}
    \definecolor{ansi-green-intense}{HTML}{007427}
    \definecolor{ansi-yellow}{HTML}{DDB62B}
    \definecolor{ansi-yellow-intense}{HTML}{B27D12}
    \definecolor{ansi-blue}{HTML}{208FFB}
    \definecolor{ansi-blue-intense}{HTML}{0065CA}
    \definecolor{ansi-magenta}{HTML}{D160C4}
    \definecolor{ansi-magenta-intense}{HTML}{A03196}
    \definecolor{ansi-cyan}{HTML}{60C6C8}
    \definecolor{ansi-cyan-intense}{HTML}{258F8F}
    \definecolor{ansi-white}{HTML}{C5C1B4}
    \definecolor{ansi-white-intense}{HTML}{A1A6B2}
    \definecolor{ansi-default-inverse-fg}{HTML}{FFFFFF}
    \definecolor{ansi-default-inverse-bg}{HTML}{000000}

    % common color for the border for error outputs.
    \definecolor{outerrorbackground}{HTML}{FFDFDF}

    % commands and environments needed by pandoc snippets
    % extracted from the output of `pandoc -s`
    \providecommand{\tightlist}{%
      \setlength{\itemsep}{0pt}\setlength{\parskip}{0pt}}
    \DefineVerbatimEnvironment{Highlighting}{Verbatim}{commandchars=\\\{\}}
    % Add ',fontsize=\small' for more characters per line
    \newenvironment{Shaded}{}{}
    \newcommand{\KeywordTok}[1]{\textcolor[rgb]{0.00,0.44,0.13}{\textbf{{#1}}}}
    \newcommand{\DataTypeTok}[1]{\textcolor[rgb]{0.56,0.13,0.00}{{#1}}}
    \newcommand{\DecValTok}[1]{\textcolor[rgb]{0.25,0.63,0.44}{{#1}}}
    \newcommand{\BaseNTok}[1]{\textcolor[rgb]{0.25,0.63,0.44}{{#1}}}
    \newcommand{\FloatTok}[1]{\textcolor[rgb]{0.25,0.63,0.44}{{#1}}}
    \newcommand{\CharTok}[1]{\textcolor[rgb]{0.25,0.44,0.63}{{#1}}}
    \newcommand{\StringTok}[1]{\textcolor[rgb]{0.25,0.44,0.63}{{#1}}}
    \newcommand{\CommentTok}[1]{\textcolor[rgb]{0.38,0.63,0.69}{\textit{{#1}}}}
    \newcommand{\OtherTok}[1]{\textcolor[rgb]{0.00,0.44,0.13}{{#1}}}
    \newcommand{\AlertTok}[1]{\textcolor[rgb]{1.00,0.00,0.00}{\textbf{{#1}}}}
    \newcommand{\FunctionTok}[1]{\textcolor[rgb]{0.02,0.16,0.49}{{#1}}}
    \newcommand{\RegionMarkerTok}[1]{{#1}}
    \newcommand{\ErrorTok}[1]{\textcolor[rgb]{1.00,0.00,0.00}{\textbf{{#1}}}}
    \newcommand{\NormalTok}[1]{{#1}}

    % Additional commands for more recent versions of Pandoc
    \newcommand{\ConstantTok}[1]{\textcolor[rgb]{0.53,0.00,0.00}{{#1}}}
    \newcommand{\SpecialCharTok}[1]{\textcolor[rgb]{0.25,0.44,0.63}{{#1}}}
    \newcommand{\VerbatimStringTok}[1]{\textcolor[rgb]{0.25,0.44,0.63}{{#1}}}
    \newcommand{\SpecialStringTok}[1]{\textcolor[rgb]{0.73,0.40,0.53}{{#1}}}
    \newcommand{\ImportTok}[1]{{#1}}
    \newcommand{\DocumentationTok}[1]{\textcolor[rgb]{0.73,0.13,0.13}{\textit{{#1}}}}
    \newcommand{\AnnotationTok}[1]{\textcolor[rgb]{0.38,0.63,0.69}{\textbf{\textit{{#1}}}}}
    \newcommand{\CommentVarTok}[1]{\textcolor[rgb]{0.38,0.63,0.69}{\textbf{\textit{{#1}}}}}
    \newcommand{\VariableTok}[1]{\textcolor[rgb]{0.10,0.09,0.49}{{#1}}}
    \newcommand{\ControlFlowTok}[1]{\textcolor[rgb]{0.00,0.44,0.13}{\textbf{{#1}}}}
    \newcommand{\OperatorTok}[1]{\textcolor[rgb]{0.40,0.40,0.40}{{#1}}}
    \newcommand{\BuiltInTok}[1]{{#1}}
    \newcommand{\ExtensionTok}[1]{{#1}}
    \newcommand{\PreprocessorTok}[1]{\textcolor[rgb]{0.74,0.48,0.00}{{#1}}}
    \newcommand{\AttributeTok}[1]{\textcolor[rgb]{0.49,0.56,0.16}{{#1}}}
    \newcommand{\InformationTok}[1]{\textcolor[rgb]{0.38,0.63,0.69}{\textbf{\textit{{#1}}}}}
    \newcommand{\WarningTok}[1]{\textcolor[rgb]{0.38,0.63,0.69}{\textbf{\textit{{#1}}}}}


    % Define a nice break command that doesn't care if a line doesn't already
    % exist.
    \def\br{\hspace*{\fill} \\* }
    % Math Jax compatibility definitions
    \def\gt{>}
    \def\lt{<}
    \let\Oldtex\TeX
    \let\Oldlatex\LaTeX
    \renewcommand{\TeX}{\textrm{\Oldtex}}
    \renewcommand{\LaTeX}{\textrm{\Oldlatex}}
    % Document parameters
    % Document title
    \title{Copia\_de\_Tuto\_1\_1\_Discovering\_Open\_Data}
    
    
    
    
    
% Pygments definitions
\makeatletter
\def\PY@reset{\let\PY@it=\relax \let\PY@bf=\relax%
    \let\PY@ul=\relax \let\PY@tc=\relax%
    \let\PY@bc=\relax \let\PY@ff=\relax}
\def\PY@tok#1{\csname PY@tok@#1\endcsname}
\def\PY@toks#1+{\ifx\relax#1\empty\else%
    \PY@tok{#1}\expandafter\PY@toks\fi}
\def\PY@do#1{\PY@bc{\PY@tc{\PY@ul{%
    \PY@it{\PY@bf{\PY@ff{#1}}}}}}}
\def\PY#1#2{\PY@reset\PY@toks#1+\relax+\PY@do{#2}}

\@namedef{PY@tok@w}{\def\PY@tc##1{\textcolor[rgb]{0.73,0.73,0.73}{##1}}}
\@namedef{PY@tok@c}{\let\PY@it=\textit\def\PY@tc##1{\textcolor[rgb]{0.24,0.48,0.48}{##1}}}
\@namedef{PY@tok@cp}{\def\PY@tc##1{\textcolor[rgb]{0.61,0.40,0.00}{##1}}}
\@namedef{PY@tok@k}{\let\PY@bf=\textbf\def\PY@tc##1{\textcolor[rgb]{0.00,0.50,0.00}{##1}}}
\@namedef{PY@tok@kp}{\def\PY@tc##1{\textcolor[rgb]{0.00,0.50,0.00}{##1}}}
\@namedef{PY@tok@kt}{\def\PY@tc##1{\textcolor[rgb]{0.69,0.00,0.25}{##1}}}
\@namedef{PY@tok@o}{\def\PY@tc##1{\textcolor[rgb]{0.40,0.40,0.40}{##1}}}
\@namedef{PY@tok@ow}{\let\PY@bf=\textbf\def\PY@tc##1{\textcolor[rgb]{0.67,0.13,1.00}{##1}}}
\@namedef{PY@tok@nb}{\def\PY@tc##1{\textcolor[rgb]{0.00,0.50,0.00}{##1}}}
\@namedef{PY@tok@nf}{\def\PY@tc##1{\textcolor[rgb]{0.00,0.00,1.00}{##1}}}
\@namedef{PY@tok@nc}{\let\PY@bf=\textbf\def\PY@tc##1{\textcolor[rgb]{0.00,0.00,1.00}{##1}}}
\@namedef{PY@tok@nn}{\let\PY@bf=\textbf\def\PY@tc##1{\textcolor[rgb]{0.00,0.00,1.00}{##1}}}
\@namedef{PY@tok@ne}{\let\PY@bf=\textbf\def\PY@tc##1{\textcolor[rgb]{0.80,0.25,0.22}{##1}}}
\@namedef{PY@tok@nv}{\def\PY@tc##1{\textcolor[rgb]{0.10,0.09,0.49}{##1}}}
\@namedef{PY@tok@no}{\def\PY@tc##1{\textcolor[rgb]{0.53,0.00,0.00}{##1}}}
\@namedef{PY@tok@nl}{\def\PY@tc##1{\textcolor[rgb]{0.46,0.46,0.00}{##1}}}
\@namedef{PY@tok@ni}{\let\PY@bf=\textbf\def\PY@tc##1{\textcolor[rgb]{0.44,0.44,0.44}{##1}}}
\@namedef{PY@tok@na}{\def\PY@tc##1{\textcolor[rgb]{0.41,0.47,0.13}{##1}}}
\@namedef{PY@tok@nt}{\let\PY@bf=\textbf\def\PY@tc##1{\textcolor[rgb]{0.00,0.50,0.00}{##1}}}
\@namedef{PY@tok@nd}{\def\PY@tc##1{\textcolor[rgb]{0.67,0.13,1.00}{##1}}}
\@namedef{PY@tok@s}{\def\PY@tc##1{\textcolor[rgb]{0.73,0.13,0.13}{##1}}}
\@namedef{PY@tok@sd}{\let\PY@it=\textit\def\PY@tc##1{\textcolor[rgb]{0.73,0.13,0.13}{##1}}}
\@namedef{PY@tok@si}{\let\PY@bf=\textbf\def\PY@tc##1{\textcolor[rgb]{0.64,0.35,0.47}{##1}}}
\@namedef{PY@tok@se}{\let\PY@bf=\textbf\def\PY@tc##1{\textcolor[rgb]{0.67,0.36,0.12}{##1}}}
\@namedef{PY@tok@sr}{\def\PY@tc##1{\textcolor[rgb]{0.64,0.35,0.47}{##1}}}
\@namedef{PY@tok@ss}{\def\PY@tc##1{\textcolor[rgb]{0.10,0.09,0.49}{##1}}}
\@namedef{PY@tok@sx}{\def\PY@tc##1{\textcolor[rgb]{0.00,0.50,0.00}{##1}}}
\@namedef{PY@tok@m}{\def\PY@tc##1{\textcolor[rgb]{0.40,0.40,0.40}{##1}}}
\@namedef{PY@tok@gh}{\let\PY@bf=\textbf\def\PY@tc##1{\textcolor[rgb]{0.00,0.00,0.50}{##1}}}
\@namedef{PY@tok@gu}{\let\PY@bf=\textbf\def\PY@tc##1{\textcolor[rgb]{0.50,0.00,0.50}{##1}}}
\@namedef{PY@tok@gd}{\def\PY@tc##1{\textcolor[rgb]{0.63,0.00,0.00}{##1}}}
\@namedef{PY@tok@gi}{\def\PY@tc##1{\textcolor[rgb]{0.00,0.52,0.00}{##1}}}
\@namedef{PY@tok@gr}{\def\PY@tc##1{\textcolor[rgb]{0.89,0.00,0.00}{##1}}}
\@namedef{PY@tok@ge}{\let\PY@it=\textit}
\@namedef{PY@tok@gs}{\let\PY@bf=\textbf}
\@namedef{PY@tok@gp}{\let\PY@bf=\textbf\def\PY@tc##1{\textcolor[rgb]{0.00,0.00,0.50}{##1}}}
\@namedef{PY@tok@go}{\def\PY@tc##1{\textcolor[rgb]{0.44,0.44,0.44}{##1}}}
\@namedef{PY@tok@gt}{\def\PY@tc##1{\textcolor[rgb]{0.00,0.27,0.87}{##1}}}
\@namedef{PY@tok@err}{\def\PY@bc##1{{\setlength{\fboxsep}{\string -\fboxrule}\fcolorbox[rgb]{1.00,0.00,0.00}{1,1,1}{\strut ##1}}}}
\@namedef{PY@tok@kc}{\let\PY@bf=\textbf\def\PY@tc##1{\textcolor[rgb]{0.00,0.50,0.00}{##1}}}
\@namedef{PY@tok@kd}{\let\PY@bf=\textbf\def\PY@tc##1{\textcolor[rgb]{0.00,0.50,0.00}{##1}}}
\@namedef{PY@tok@kn}{\let\PY@bf=\textbf\def\PY@tc##1{\textcolor[rgb]{0.00,0.50,0.00}{##1}}}
\@namedef{PY@tok@kr}{\let\PY@bf=\textbf\def\PY@tc##1{\textcolor[rgb]{0.00,0.50,0.00}{##1}}}
\@namedef{PY@tok@bp}{\def\PY@tc##1{\textcolor[rgb]{0.00,0.50,0.00}{##1}}}
\@namedef{PY@tok@fm}{\def\PY@tc##1{\textcolor[rgb]{0.00,0.00,1.00}{##1}}}
\@namedef{PY@tok@vc}{\def\PY@tc##1{\textcolor[rgb]{0.10,0.09,0.49}{##1}}}
\@namedef{PY@tok@vg}{\def\PY@tc##1{\textcolor[rgb]{0.10,0.09,0.49}{##1}}}
\@namedef{PY@tok@vi}{\def\PY@tc##1{\textcolor[rgb]{0.10,0.09,0.49}{##1}}}
\@namedef{PY@tok@vm}{\def\PY@tc##1{\textcolor[rgb]{0.10,0.09,0.49}{##1}}}
\@namedef{PY@tok@sa}{\def\PY@tc##1{\textcolor[rgb]{0.73,0.13,0.13}{##1}}}
\@namedef{PY@tok@sb}{\def\PY@tc##1{\textcolor[rgb]{0.73,0.13,0.13}{##1}}}
\@namedef{PY@tok@sc}{\def\PY@tc##1{\textcolor[rgb]{0.73,0.13,0.13}{##1}}}
\@namedef{PY@tok@dl}{\def\PY@tc##1{\textcolor[rgb]{0.73,0.13,0.13}{##1}}}
\@namedef{PY@tok@s2}{\def\PY@tc##1{\textcolor[rgb]{0.73,0.13,0.13}{##1}}}
\@namedef{PY@tok@sh}{\def\PY@tc##1{\textcolor[rgb]{0.73,0.13,0.13}{##1}}}
\@namedef{PY@tok@s1}{\def\PY@tc##1{\textcolor[rgb]{0.73,0.13,0.13}{##1}}}
\@namedef{PY@tok@mb}{\def\PY@tc##1{\textcolor[rgb]{0.40,0.40,0.40}{##1}}}
\@namedef{PY@tok@mf}{\def\PY@tc##1{\textcolor[rgb]{0.40,0.40,0.40}{##1}}}
\@namedef{PY@tok@mh}{\def\PY@tc##1{\textcolor[rgb]{0.40,0.40,0.40}{##1}}}
\@namedef{PY@tok@mi}{\def\PY@tc##1{\textcolor[rgb]{0.40,0.40,0.40}{##1}}}
\@namedef{PY@tok@il}{\def\PY@tc##1{\textcolor[rgb]{0.40,0.40,0.40}{##1}}}
\@namedef{PY@tok@mo}{\def\PY@tc##1{\textcolor[rgb]{0.40,0.40,0.40}{##1}}}
\@namedef{PY@tok@ch}{\let\PY@it=\textit\def\PY@tc##1{\textcolor[rgb]{0.24,0.48,0.48}{##1}}}
\@namedef{PY@tok@cm}{\let\PY@it=\textit\def\PY@tc##1{\textcolor[rgb]{0.24,0.48,0.48}{##1}}}
\@namedef{PY@tok@cpf}{\let\PY@it=\textit\def\PY@tc##1{\textcolor[rgb]{0.24,0.48,0.48}{##1}}}
\@namedef{PY@tok@c1}{\let\PY@it=\textit\def\PY@tc##1{\textcolor[rgb]{0.24,0.48,0.48}{##1}}}
\@namedef{PY@tok@cs}{\let\PY@it=\textit\def\PY@tc##1{\textcolor[rgb]{0.24,0.48,0.48}{##1}}}

\def\PYZbs{\char`\\}
\def\PYZus{\char`\_}
\def\PYZob{\char`\{}
\def\PYZcb{\char`\}}
\def\PYZca{\char`\^}
\def\PYZam{\char`\&}
\def\PYZlt{\char`\<}
\def\PYZgt{\char`\>}
\def\PYZsh{\char`\#}
\def\PYZpc{\char`\%}
\def\PYZdl{\char`\$}
\def\PYZhy{\char`\-}
\def\PYZsq{\char`\'}
\def\PYZdq{\char`\"}
\def\PYZti{\char`\~}
% for compatibility with earlier versions
\def\PYZat{@}
\def\PYZlb{[}
\def\PYZrb{]}
\makeatother


    % For linebreaks inside Verbatim environment from package fancyvrb.
    \makeatletter
        \newbox\Wrappedcontinuationbox
        \newbox\Wrappedvisiblespacebox
        \newcommand*\Wrappedvisiblespace {\textcolor{red}{\textvisiblespace}}
        \newcommand*\Wrappedcontinuationsymbol {\textcolor{red}{\llap{\tiny$\m@th\hookrightarrow$}}}
        \newcommand*\Wrappedcontinuationindent {3ex }
        \newcommand*\Wrappedafterbreak {\kern\Wrappedcontinuationindent\copy\Wrappedcontinuationbox}
        % Take advantage of the already applied Pygments mark-up to insert
        % potential linebreaks for TeX processing.
        %        {, <, #, %, $, ' and ": go to next line.
        %        _, }, ^, &, >, - and ~: stay at end of broken line.
        % Use of \textquotesingle for straight quote.
        \newcommand*\Wrappedbreaksatspecials {%
            \def\PYGZus{\discretionary{\char`\_}{\Wrappedafterbreak}{\char`\_}}%
            \def\PYGZob{\discretionary{}{\Wrappedafterbreak\char`\{}{\char`\{}}%
            \def\PYGZcb{\discretionary{\char`\}}{\Wrappedafterbreak}{\char`\}}}%
            \def\PYGZca{\discretionary{\char`\^}{\Wrappedafterbreak}{\char`\^}}%
            \def\PYGZam{\discretionary{\char`\&}{\Wrappedafterbreak}{\char`\&}}%
            \def\PYGZlt{\discretionary{}{\Wrappedafterbreak\char`\<}{\char`\<}}%
            \def\PYGZgt{\discretionary{\char`\>}{\Wrappedafterbreak}{\char`\>}}%
            \def\PYGZsh{\discretionary{}{\Wrappedafterbreak\char`\#}{\char`\#}}%
            \def\PYGZpc{\discretionary{}{\Wrappedafterbreak\char`\%}{\char`\%}}%
            \def\PYGZdl{\discretionary{}{\Wrappedafterbreak\char`\$}{\char`\$}}%
            \def\PYGZhy{\discretionary{\char`\-}{\Wrappedafterbreak}{\char`\-}}%
            \def\PYGZsq{\discretionary{}{\Wrappedafterbreak\textquotesingle}{\textquotesingle}}%
            \def\PYGZdq{\discretionary{}{\Wrappedafterbreak\char`\"}{\char`\"}}%
            \def\PYGZti{\discretionary{\char`\~}{\Wrappedafterbreak}{\char`\~}}%
        }
        % Some characters . , ; ? ! / are not pygmentized.
        % This macro makes them "active" and they will insert potential linebreaks
        \newcommand*\Wrappedbreaksatpunct {%
            \lccode`\~`\.\lowercase{\def~}{\discretionary{\hbox{\char`\.}}{\Wrappedafterbreak}{\hbox{\char`\.}}}%
            \lccode`\~`\,\lowercase{\def~}{\discretionary{\hbox{\char`\,}}{\Wrappedafterbreak}{\hbox{\char`\,}}}%
            \lccode`\~`\;\lowercase{\def~}{\discretionary{\hbox{\char`\;}}{\Wrappedafterbreak}{\hbox{\char`\;}}}%
            \lccode`\~`\:\lowercase{\def~}{\discretionary{\hbox{\char`\:}}{\Wrappedafterbreak}{\hbox{\char`\:}}}%
            \lccode`\~`\?\lowercase{\def~}{\discretionary{\hbox{\char`\?}}{\Wrappedafterbreak}{\hbox{\char`\?}}}%
            \lccode`\~`\!\lowercase{\def~}{\discretionary{\hbox{\char`\!}}{\Wrappedafterbreak}{\hbox{\char`\!}}}%
            \lccode`\~`\/\lowercase{\def~}{\discretionary{\hbox{\char`\/}}{\Wrappedafterbreak}{\hbox{\char`\/}}}%
            \catcode`\.\active
            \catcode`\,\active
            \catcode`\;\active
            \catcode`\:\active
            \catcode`\?\active
            \catcode`\!\active
            \catcode`\/\active
            \lccode`\~`\~
        }
    \makeatother

    \let\OriginalVerbatim=\Verbatim
    \makeatletter
    \renewcommand{\Verbatim}[1][1]{%
        %\parskip\z@skip
        \sbox\Wrappedcontinuationbox {\Wrappedcontinuationsymbol}%
        \sbox\Wrappedvisiblespacebox {\FV@SetupFont\Wrappedvisiblespace}%
        \def\FancyVerbFormatLine ##1{\hsize\linewidth
            \vtop{\raggedright\hyphenpenalty\z@\exhyphenpenalty\z@
                \doublehyphendemerits\z@\finalhyphendemerits\z@
                \strut ##1\strut}%
        }%
        % If the linebreak is at a space, the latter will be displayed as visible
        % space at end of first line, and a continuation symbol starts next line.
        % Stretch/shrink are however usually zero for typewriter font.
        \def\FV@Space {%
            \nobreak\hskip\z@ plus\fontdimen3\font minus\fontdimen4\font
            \discretionary{\copy\Wrappedvisiblespacebox}{\Wrappedafterbreak}
            {\kern\fontdimen2\font}%
        }%

        % Allow breaks at special characters using \PYG... macros.
        \Wrappedbreaksatspecials
        % Breaks at punctuation characters . , ; ? ! and / need catcode=\active
        \OriginalVerbatim[#1,codes*=\Wrappedbreaksatpunct]%
    }
    \makeatother

    % Exact colors from NB
    \definecolor{incolor}{HTML}{303F9F}
    \definecolor{outcolor}{HTML}{D84315}
    \definecolor{cellborder}{HTML}{CFCFCF}
    \definecolor{cellbackground}{HTML}{F7F7F7}

    % prompt
    \makeatletter
    \newcommand{\boxspacing}{\kern\kvtcb@left@rule\kern\kvtcb@boxsep}
    \makeatother
    \newcommand{\prompt}[4]{
        {\ttfamily\llap{{\color{#2}[#3]:\hspace{3pt}#4}}\vspace{-\baselineskip}}
    }
    

    
    % Prevent overflowing lines due to hard-to-break entities
    \sloppy
    % Setup hyperref package
    \hypersetup{
      breaklinks=true,  % so long urls are correctly broken across lines
      colorlinks=true,
      urlcolor=urlcolor,
      linkcolor=linkcolor,
      citecolor=citecolor,
      }
    % Slightly bigger margins than the latex defaults
    
    \geometry{verbose,tmargin=1in,bmargin=1in,lmargin=1in,rmargin=1in}
    
    

\begin{document}
    
    \maketitle
    
    

    
    \hypertarget{gravitational-wave-open-data-workshop-3}{%
\section{Gravitational Wave Open Data Workshop
\#3}\label{gravitational-wave-open-data-workshop-3}}

\hypertarget{tutorial-1.1-discovering-open-data-from-gw-observatories}{%
\paragraph{Tutorial 1.1: Discovering open data from GW
observatories}\label{tutorial-1.1-discovering-open-data-from-gw-observatories}}

This notebook describes how to discover what data are available from the
\href{https://www.gw-openscience.org}{Gravitational-Wave Open Science
Center (GWOSC)}.

\href{https://colab.research.google.com/github/gw-odw/odw-2020/blob/master/Day_1/Tuto\%201.1\%20Discovering\%20Open\%20Data.ipynb}{Click
this link to view this tutorial in Google Colaboratory}

    \hypertarget{software-installation-execute-only-if-running-on-a-cloud-platform-or-havent-done-the-installation-yet}{%
\subsection{Software installation (execute only if running on a cloud
platform or haven't done the installation
yet!)}\label{software-installation-execute-only-if-running-on-a-cloud-platform-or-havent-done-the-installation-yet}}

First, we need to install the software, which we do following the
instruction in
\href{https://github.com/gw-odw/odw-2020/blob/master/setup.md}{Software
Setup Instructions}:

    \begin{tcolorbox}[breakable, size=fbox, boxrule=1pt, pad at break*=1mm,colback=cellbackground, colframe=cellborder]
\prompt{In}{incolor}{1}{\boxspacing}
\begin{Verbatim}[commandchars=\\\{\}]
\PY{c+c1}{\PYZsh{} \PYZhy{}\PYZhy{} Uncomment following line if running in Google Colab}
\PY{o}{!}pip\PY{+w}{ }install\PY{+w}{ }\PYZhy{}U\PY{+w}{ }gwosc
\end{Verbatim}
\end{tcolorbox}

    \begin{Verbatim}[commandchars=\\\{\}]
\textcolor{ansi-yellow}{WARNING: Skipping
/home/debian12/Documentos/UD/Mineria\_de\_Datos/MD/lib/python3.11/site-
packages/nltk-3.9.2.dist-info due to invalid metadata entry 'name'}\textcolor{ansi-yellow}{
}\textcolor{ansi-yellow}{WARNING: Skipping
/home/debian12/Documentos/UD/Mineria\_de\_Datos/MD/lib/python3.11/site-
packages/wordcloud-1.9.4.dist-info due to invalid metadata entry 'name'}\textcolor{ansi-yellow}{
}\textcolor{ansi-yellow}{WARNING: Skipping
/home/debian12/Documentos/UD/Mineria\_de\_Datos/MD/lib/python3.11/site-
packages/nltk-3.9.2.dist-info due to invalid metadata entry 'name'}\textcolor{ansi-yellow}{
}\textcolor{ansi-yellow}{WARNING: Skipping
/home/debian12/Documentos/UD/Mineria\_de\_Datos/MD/lib/python3.11/site-
packages/wordcloud-1.9.4.dist-info due to invalid metadata entry 'name'}\textcolor{ansi-yellow}{
}\textcolor{ansi-yellow}{WARNING: Skipping
/home/debian12/Documentos/UD/Mineria\_de\_Datos/MD/lib/python3.11/site-
packages/regex-2025.11.3.dist-info due to invalid metadata entry 'name'}\textcolor{ansi-yellow}{
}Collecting gwosc
  Downloading gwosc-0.8.1-py3-none-any.whl (32 kB)
Requirement already satisfied: requests>=1.0.0 in
/home/debian12/Documentos/UD/Mineria\_de\_Datos/MD/lib/python3.11/site-packages
(from gwosc) (2.32.5)
Requirement already satisfied: charset\_normalizer<4,>=2 in
/home/debian12/Documentos/UD/Mineria\_de\_Datos/MD/lib/python3.11/site-packages
(from requests>=1.0.0->gwosc) (3.4.4)
Requirement already satisfied: idna<4,>=2.5 in
/home/debian12/Documentos/UD/Mineria\_de\_Datos/MD/lib/python3.11/site-packages
(from requests>=1.0.0->gwosc) (3.11)
Requirement already satisfied: urllib3<3,>=1.21.1 in
/home/debian12/Documentos/UD/Mineria\_de\_Datos/MD/lib/python3.11/site-packages
(from requests>=1.0.0->gwosc) (2.5.0)
Requirement already satisfied: certifi>=2017.4.17 in
/home/debian12/Documentos/UD/Mineria\_de\_Datos/MD/lib/python3.11/site-packages
(from requests>=1.0.0->gwosc) (2025.11.12)
\textcolor{ansi-yellow}{WARNING: Skipping
/home/debian12/Documentos/UD/Mineria\_de\_Datos/MD/lib/python3.11/site-
packages/nltk-3.9.2.dist-info due to invalid metadata entry 'name'}\textcolor{ansi-yellow}{
}\textcolor{ansi-yellow}{WARNING: Skipping
/home/debian12/Documentos/UD/Mineria\_de\_Datos/MD/lib/python3.11/site-
packages/wordcloud-1.9.4.dist-info due to invalid metadata entry 'name'}\textcolor{ansi-yellow}{
}\textcolor{ansi-yellow}{WARNING: Skipping
/home/debian12/Documentos/UD/Mineria\_de\_Datos/MD/lib/python3.11/site-
packages/regex-2025.11.3.dist-info due to invalid metadata entry 'name'}\textcolor{ansi-yellow}{
}Installing collected packages: gwosc
\textcolor{ansi-yellow}{WARNING: Skipping
/home/debian12/Documentos/UD/Mineria\_de\_Datos/MD/lib/python3.11/site-
packages/nltk-3.9.2.dist-info due to invalid metadata entry 'name'}\textcolor{ansi-yellow}{
}\textcolor{ansi-yellow}{WARNING: Skipping
/home/debian12/Documentos/UD/Mineria\_de\_Datos/MD/lib/python3.11/site-
packages/wordcloud-1.9.4.dist-info due to invalid metadata entry 'name'}\textcolor{ansi-yellow}{
}\textcolor{ansi-yellow}{WARNING: Skipping
/home/debian12/Documentos/UD/Mineria\_de\_Datos/MD/lib/python3.11/site-
packages/regex-2025.11.3.dist-info due to invalid metadata entry 'name'}\textcolor{ansi-yellow}{
}Successfully installed gwosc-0.8.1
    \end{Verbatim}

    \textbf{Important:} With Google Colab, you may need to restart the
runtime after running the cell above.

    \begin{tcolorbox}[breakable, size=fbox, boxrule=1pt, pad at break*=1mm,colback=cellbackground, colframe=cellborder]
\prompt{In}{incolor}{2}{\boxspacing}
\begin{Verbatim}[commandchars=\\\{\}]
\PY{c+c1}{\PYZsh{}check the version of the package gwosc you are using}
\PY{k+kn}{import} \PY{n+nn}{gwosc}
\PY{n+nb}{print}\PY{p}{(}\PY{n}{gwosc}\PY{o}{.}\PY{n}{\PYZus{}\PYZus{}version\PYZus{}\PYZus{}}\PY{p}{)}
\end{Verbatim}
\end{tcolorbox}

    \begin{Verbatim}[commandchars=\\\{\}]
0.8.1
    \end{Verbatim}

    The version you get should be 0.8 or greater. If it's not, check that
you have followed all the steps in
\href{https://github.com/gw-odw/odw-2020/blob/master/setup.md}{Software
Setup Instructions}.

    \hypertarget{querying-for-event-information}{%
\subsection{Querying for event
information}\label{querying-for-event-information}}

The module \texttt{gwosc.catalog} provides tools to search for events in
a catalog.

The module \texttt{gwosc.datasets} provides tools for searching for
datasets, including full run strain data releases.

For example, we can search for events in the
\href{https://www.gw-openscience.org/eventapi/html/GWTC-1-confident/}{GWTC-1
catalog}, the catalog of all events from the O1 and O2 observing runs. A
list of available catalogs can be seen in the
\href{https://gw-openscience.org/eventapi}{Event Portal}

    \begin{tcolorbox}[breakable, size=fbox, boxrule=1pt, pad at break*=1mm,colback=cellbackground, colframe=cellborder]
\prompt{In}{incolor}{3}{\boxspacing}
\begin{Verbatim}[commandchars=\\\{\}]
\PY{k+kn}{from} \PY{n+nn}{gwosc}\PY{n+nn}{.}\PY{n+nn}{datasets} \PY{k+kn}{import} \PY{n}{find\PYZus{}datasets}

\PY{c+c1}{\PYZsh{} \PYZhy{}\PYZhy{} GW events del catálogo GWTC\PYZhy{}1\PYZhy{}confident}
\PY{n}{gwtc1} \PY{o}{=} \PY{n}{find\PYZus{}datasets}\PY{p}{(}\PY{n+nb}{type}\PY{o}{=}\PY{l+s+s2}{\PYZdq{}}\PY{l+s+s2}{event}\PY{l+s+s2}{\PYZdq{}}\PY{p}{,} \PY{n}{catalog}\PY{o}{=}\PY{l+s+s2}{\PYZdq{}}\PY{l+s+s2}{GWTC\PYZhy{}1\PYZhy{}confident}\PY{l+s+s2}{\PYZdq{}}\PY{p}{)}
\PY{n+nb}{print}\PY{p}{(}\PY{l+s+s2}{\PYZdq{}}\PY{l+s+s2}{GWTC\PYZhy{}1 events:}\PY{l+s+s2}{\PYZdq{}}\PY{p}{,} \PY{n}{gwtc1}\PY{p}{)}
\PY{n+nb}{print}\PY{p}{(}\PY{l+s+s2}{\PYZdq{}}\PY{l+s+s2}{\PYZdq{}}\PY{p}{)}

\PY{c+c1}{\PYZsh{} \PYZhy{}\PYZhy{} Conjuntos grandes de datos de strain (runs de LIGO/Virgo)}
\PY{n}{runs} \PY{o}{=} \PY{n}{find\PYZus{}datasets}\PY{p}{(}\PY{n+nb}{type}\PY{o}{=}\PY{l+s+s2}{\PYZdq{}}\PY{l+s+s2}{run}\PY{l+s+s2}{\PYZdq{}}\PY{p}{)}
\PY{n+nb}{print}\PY{p}{(}\PY{l+s+s2}{\PYZdq{}}\PY{l+s+s2}{Large data sets:}\PY{l+s+s2}{\PYZdq{}}\PY{p}{,} \PY{n}{runs}\PY{p}{)}
\end{Verbatim}
\end{tcolorbox}

    \begin{Verbatim}[commandchars=\\\{\}]
GWTC-1 events: ['GW150914-v3', 'GW151012-v3', 'GW151226-v2', 'GW170104-v2',
'GW170608-v3', 'GW170729-v1', 'GW170809-v1', 'GW170814-v3', 'GW170817-v3',
'GW170818-v1', 'GW170823-v1']

Large data sets: ['O1', 'O2', 'O3GK', 'O3a', 'O3b', 'O4a', 'O4b1Disc',
'O4b2Disc', 'O4b3Disc', 'S5', 'S6']
    \end{Verbatim}

    \texttt{catalog.events} also accepts a \texttt{segment} and
\texttt{detector} keyword to narrow results based on GPS time and
detector:

    \begin{tcolorbox}[breakable, size=fbox, boxrule=1pt, pad at break*=1mm,colback=cellbackground, colframe=cellborder]
\prompt{In}{incolor}{4}{\boxspacing}
\begin{Verbatim}[commandchars=\\\{\}]
\PY{c+c1}{\PYZsh{}\PYZhy{}\PYZhy{} Detector and segments keywords limit search results}

\PY{n}{datasets} \PY{o}{=} \PY{n}{find\PYZus{}datasets}\PY{p}{(}
    \PY{n+nb}{type}\PY{o}{=}\PY{l+s+s2}{\PYZdq{}}\PY{l+s+s2}{event}\PY{l+s+s2}{\PYZdq{}}\PY{p}{,}
    \PY{n}{catalog}\PY{o}{=}\PY{l+s+s2}{\PYZdq{}}\PY{l+s+s2}{GWTC\PYZhy{}1\PYZhy{}confident}\PY{l+s+s2}{\PYZdq{}}\PY{p}{,}
    \PY{n}{detector}\PY{o}{=}\PY{l+s+s2}{\PYZdq{}}\PY{l+s+s2}{L1}\PY{l+s+s2}{\PYZdq{}}\PY{p}{,}
    \PY{n}{segment}\PY{o}{=}\PY{p}{(}\PY{l+m+mi}{1164556817}\PY{p}{,} \PY{l+m+mi}{1187733618}\PY{p}{)}
\PY{p}{)}

\PY{n+nb}{print}\PY{p}{(}\PY{n}{datasets}\PY{p}{)}
\end{Verbatim}
\end{tcolorbox}

    \begin{Verbatim}[commandchars=\\\{\}]
['GW170104-v2', 'GW170608-v3', 'GW170729-v1', 'GW170809-v1', 'GW170814-v3',
'GW170817-v3', 'GW170818-v1', 'GW170823-v1']
    \end{Verbatim}

    Using \texttt{gwosc.datasets.event\_gps}, we can query for the GPS time
of a specific event:

    \begin{tcolorbox}[breakable, size=fbox, boxrule=1pt, pad at break*=1mm,colback=cellbackground, colframe=cellborder]
\prompt{In}{incolor}{5}{\boxspacing}
\begin{Verbatim}[commandchars=\\\{\}]
\PY{k+kn}{from} \PY{n+nn}{gwosc}\PY{n+nn}{.}\PY{n+nn}{datasets} \PY{k+kn}{import} \PY{n}{event\PYZus{}gps}
\PY{n}{gps} \PY{o}{=} \PY{n}{event\PYZus{}gps}\PY{p}{(}\PY{l+s+s1}{\PYZsq{}}\PY{l+s+s1}{GW190425}\PY{l+s+s1}{\PYZsq{}}\PY{p}{)}
\PY{n+nb}{print}\PY{p}{(}\PY{n}{gps}\PY{p}{)}
\end{Verbatim}
\end{tcolorbox}

    \begin{Verbatim}[commandchars=\\\{\}]
1240215503.0
    \end{Verbatim}

    All of these times are returned in the GPS time system, which counts the
number of seconds that have elapsed since the start of the GPS epoch at
midnight (00:00) on January 6th 1980. GWOSC provides a GPS time
converter you can use to translate into datetime, or you can use
gwpy.time.

    We can query for the GPS time interval for an observing run:

    \begin{tcolorbox}[breakable, size=fbox, boxrule=1pt, pad at break*=1mm,colback=cellbackground, colframe=cellborder]
\prompt{In}{incolor}{7}{\boxspacing}
\begin{Verbatim}[commandchars=\\\{\}]
\PY{k+kn}{from} \PY{n+nn}{gwosc}\PY{n+nn}{.}\PY{n+nn}{datasets} \PY{k+kn}{import} \PY{n}{run\PYZus{}segment}
\PY{n+nb}{print}\PY{p}{(}\PY{n}{run\PYZus{}segment}\PY{p}{(}\PY{l+s+s1}{\PYZsq{}}\PY{l+s+s1}{O1}\PY{l+s+s1}{\PYZsq{}}\PY{p}{)}\PY{p}{)}
\end{Verbatim}
\end{tcolorbox}

    \begin{Verbatim}[commandchars=\\\{\}]
(1126051217, 1137254417)
    \end{Verbatim}

    To see only the confident events in O1:

    \begin{tcolorbox}[breakable, size=fbox, boxrule=1pt, pad at break*=1mm,colback=cellbackground, colframe=cellborder]
\prompt{In}{incolor}{8}{\boxspacing}
\begin{Verbatim}[commandchars=\\\{\}]
\PY{c+c1}{\PYZsh{} Segmento GPS del run O1 (2015–2016)}
\PY{n}{O1\PYZus{}segment} \PY{o}{=} \PY{p}{(}\PY{l+m+mi}{1126051217}\PY{p}{,} \PY{l+m+mi}{1137254417}\PY{p}{)}

\PY{c+c1}{\PYZsh{} Buscar los eventos del catálogo GWTC\PYZhy{}1\PYZhy{}confident que ocurrieron en O1}
\PY{n}{O1\PYZus{}events} \PY{o}{=} \PY{n}{find\PYZus{}datasets}\PY{p}{(}
    \PY{n+nb}{type}\PY{o}{=}\PY{l+s+s2}{\PYZdq{}}\PY{l+s+s2}{event}\PY{l+s+s2}{\PYZdq{}}\PY{p}{,}
    \PY{n}{catalog}\PY{o}{=}\PY{l+s+s2}{\PYZdq{}}\PY{l+s+s2}{GWTC\PYZhy{}1\PYZhy{}confident}\PY{l+s+s2}{\PYZdq{}}\PY{p}{,}
    \PY{n}{segment}\PY{o}{=}\PY{n}{O1\PYZus{}segment}
\PY{p}{)}

\PY{n+nb}{print}\PY{p}{(}\PY{l+s+s2}{\PYZdq{}}\PY{l+s+s2}{Eventos en O1:}\PY{l+s+s2}{\PYZdq{}}\PY{p}{,} \PY{n}{O1\PYZus{}events}\PY{p}{)}
\end{Verbatim}
\end{tcolorbox}

    \begin{Verbatim}[commandchars=\\\{\}]
Eventos en O1: ['GW150914-v3', 'GW151012-v3', 'GW151226-v2']
    \end{Verbatim}

    \hypertarget{querying-for-data-files}{%
\subsection{Querying for data files}\label{querying-for-data-files}}

The \texttt{gwosc.locate} module provides a function to find the URLs of
data files associated with a given dataset.

For event datasets, one can get the list of URLs using only the event
name:

    \begin{tcolorbox}[breakable, size=fbox, boxrule=1pt, pad at break*=1mm,colback=cellbackground, colframe=cellborder]
\prompt{In}{incolor}{9}{\boxspacing}
\begin{Verbatim}[commandchars=\\\{\}]
\PY{k+kn}{from} \PY{n+nn}{gwosc}\PY{n+nn}{.}\PY{n+nn}{locate} \PY{k+kn}{import} \PY{n}{get\PYZus{}event\PYZus{}urls}
\PY{n}{urls} \PY{o}{=} \PY{n}{get\PYZus{}event\PYZus{}urls}\PY{p}{(}\PY{l+s+s1}{\PYZsq{}}\PY{l+s+s1}{GW150914}\PY{l+s+s1}{\PYZsq{}}\PY{p}{)}
\PY{n+nb}{print}\PY{p}{(}\PY{n}{urls}\PY{p}{)}
\end{Verbatim}
\end{tcolorbox}

    \begin{Verbatim}[commandchars=\\\{\}]
['https://gwosc.org/archive/data/O1/1126170624/H-H1\_LOSC\_4\_V1-1126256640-
4096.hdf5', 'https://gwosc.org/archive/data/O1/1126170624/L-L1\_LOSC\_4\_V1-
1126256640-4096.hdf5']
    \end{Verbatim}

    By default, this function returns all of the files associated with a
given event, which isn't particularly helpful. However, we can can
filter on any of these by using keyword arguments, for example to get
the URL for the 32-second file for the LIGO-Livingston detector:

    \begin{tcolorbox}[breakable, size=fbox, boxrule=1pt, pad at break*=1mm,colback=cellbackground, colframe=cellborder]
\prompt{In}{incolor}{10}{\boxspacing}
\begin{Verbatim}[commandchars=\\\{\}]
\PY{n}{urls} \PY{o}{=} \PY{n}{get\PYZus{}event\PYZus{}urls}\PY{p}{(}\PY{l+s+s1}{\PYZsq{}}\PY{l+s+s1}{GW150914}\PY{l+s+s1}{\PYZsq{}}\PY{p}{,} \PY{n}{duration}\PY{o}{=}\PY{l+m+mi}{32}\PY{p}{,} \PY{n}{detector}\PY{o}{=}\PY{l+s+s1}{\PYZsq{}}\PY{l+s+s1}{L1}\PY{l+s+s1}{\PYZsq{}}\PY{p}{)}
\PY{n+nb}{print}\PY{p}{(}\PY{n}{urls}\PY{p}{)}
\end{Verbatim}
\end{tcolorbox}

    \begin{Verbatim}[commandchars=\\\{\}]
['https://gwosc.org/archive/data/O1/1126170624/L-L1\_LOSC\_4\_V1-1126256640-
4096.hdf5']
    \end{Verbatim}

    \hypertarget{exercises}{%
\section{Exercises}\label{exercises}}

Ahora que has visto ejemplos de cómo consultar información sobre
conjuntos de datos usando el paquete \texttt{gwosc}, intenta completar
los siguientes ejercicios utilizando esa interfaz:

\begin{itemize}
\item
  ¿Cuántos meses duró S6?
\item
  ¿Cuántos eventos del catálogo GWTC-1-confident fueron detectados
  durante O1?
\item
  ¿Qué URL de archivo contiene los datos para el detector V1 en un
  intervalo de 4096 segundos alrededor de GW170817?
\end{itemize}

    \hypertarget{soluciuxf3n}{%
\section{SOLUCIÓN}\label{soluciuxf3n}}

    \hypertarget{cuuxe1ntos-meses-duruxf3-s6}{%
\subsection{¿Cuántos meses duró S6?}\label{cuuxe1ntos-meses-duruxf3-s6}}

    \begin{tcolorbox}[breakable, size=fbox, boxrule=1pt, pad at break*=1mm,colback=cellbackground, colframe=cellborder]
\prompt{In}{incolor}{13}{\boxspacing}
\begin{Verbatim}[commandchars=\\\{\}]
\PY{k+kn}{from} \PY{n+nn}{gwosc}\PY{n+nn}{.}\PY{n+nn}{datasets} \PY{k+kn}{import} \PY{n}{run\PYZus{}segment}
\PY{n+nb}{print}\PY{p}{(}\PY{n}{run\PYZus{}segment}\PY{p}{(}\PY{l+s+s1}{\PYZsq{}}\PY{l+s+s1}{S6}\PY{l+s+s1}{\PYZsq{}}\PY{p}{)}\PY{p}{)}
\end{Verbatim}
\end{tcolorbox}

    \begin{Verbatim}[commandchars=\\\{\}]
(931035615, 971622015)
    \end{Verbatim}

    Al realizar la conversión del tiempo GPS al tiempo del calendario
gregoriano se tiene:

\begin{itemize}
\tightlist
\item
  931035615 = 2009-07-07 T 21:00:00
\item
  971622015 = 2010-10-20 T 15:00:00
\end{itemize}

Por lo tanto S6 duró 1 año y 3 meses

    \hypertarget{cuuxe1ntos-eventos-del-catuxe1logo-gwtc-1-confident-fueron-detectados-durante-o1}{%
\subsection{¿Cuántos eventos del catálogo GWTC-1-confident fueron
detectados durante
O1?}\label{cuuxe1ntos-eventos-del-catuxe1logo-gwtc-1-confident-fueron-detectados-durante-o1}}

    \begin{tcolorbox}[breakable, size=fbox, boxrule=1pt, pad at break*=1mm,colback=cellbackground, colframe=cellborder]
\prompt{In}{incolor}{14}{\boxspacing}
\begin{Verbatim}[commandchars=\\\{\}]
\PY{k+kn}{from} \PY{n+nn}{gwosc}\PY{n+nn}{.}\PY{n+nn}{datasets} \PY{k+kn}{import} \PY{n}{run\PYZus{}segment}
\PY{n+nb}{print}\PY{p}{(}\PY{n}{run\PYZus{}segment}\PY{p}{(}\PY{l+s+s1}{\PYZsq{}}\PY{l+s+s1}{O1}\PY{l+s+s1}{\PYZsq{}}\PY{p}{)}\PY{p}{)}
\end{Verbatim}
\end{tcolorbox}

    \begin{Verbatim}[commandchars=\\\{\}]
(1126051217, 1137254417)
    \end{Verbatim}

    \begin{tcolorbox}[breakable, size=fbox, boxrule=1pt, pad at break*=1mm,colback=cellbackground, colframe=cellborder]
\prompt{In}{incolor}{15}{\boxspacing}
\begin{Verbatim}[commandchars=\\\{\}]
\PY{c+c1}{\PYZsh{} Segmento GPS del run O1 (2015–2016)}
\PY{n}{O1\PYZus{}segment} \PY{o}{=} \PY{p}{(}\PY{l+m+mi}{1126051217}\PY{p}{,} \PY{l+m+mi}{1137254417}\PY{p}{)}

\PY{c+c1}{\PYZsh{} Buscar los eventos del catálogo GWTC\PYZhy{}1\PYZhy{}confident que ocurrieron en O1}
\PY{n}{O1\PYZus{}events} \PY{o}{=} \PY{n}{find\PYZus{}datasets}\PY{p}{(}
    \PY{n+nb}{type}\PY{o}{=}\PY{l+s+s2}{\PYZdq{}}\PY{l+s+s2}{event}\PY{l+s+s2}{\PYZdq{}}\PY{p}{,}
    \PY{n}{catalog}\PY{o}{=}\PY{l+s+s2}{\PYZdq{}}\PY{l+s+s2}{GWTC\PYZhy{}1\PYZhy{}confident}\PY{l+s+s2}{\PYZdq{}}\PY{p}{,}
    \PY{n}{segment}\PY{o}{=}\PY{n}{O1\PYZus{}segment}
\PY{p}{)}

\PY{n+nb}{print}\PY{p}{(}\PY{l+s+s2}{\PYZdq{}}\PY{l+s+s2}{Eventos en O1:}\PY{l+s+s2}{\PYZdq{}}\PY{p}{,} \PY{n}{O1\PYZus{}events}\PY{p}{)}
\end{Verbatim}
\end{tcolorbox}

    \begin{Verbatim}[commandchars=\\\{\}]
Eventos en O1: ['GW150914-v3', 'GW151012-v3', 'GW151226-v2']
    \end{Verbatim}

    Solo hay 3 eventos registrados para ese segmento de tiempo

    \hypertarget{quuxe9-url-de-archivo-contiene-los-datos-para-el-detector-v1-en-un-intervalo-de-4096-segundos-alrededor-de-gw170817}{%
\subsection{¿Qué URL de archivo contiene los datos para el detector V1
en un intervalo de 4096 segundos alrededor de
GW170817?}\label{quuxe9-url-de-archivo-contiene-los-datos-para-el-detector-v1-en-un-intervalo-de-4096-segundos-alrededor-de-gw170817}}

    \begin{tcolorbox}[breakable, size=fbox, boxrule=1pt, pad at break*=1mm,colback=cellbackground, colframe=cellborder]
\prompt{In}{incolor}{16}{\boxspacing}
\begin{Verbatim}[commandchars=\\\{\}]
\PY{k+kn}{from} \PY{n+nn}{gwosc}\PY{n+nn}{.}\PY{n+nn}{locate} \PY{k+kn}{import} \PY{n}{get\PYZus{}event\PYZus{}urls}
\PY{n}{urls} \PY{o}{=} \PY{n}{get\PYZus{}event\PYZus{}urls}\PY{p}{(}\PY{l+s+s1}{\PYZsq{}}\PY{l+s+s1}{GW170817}\PY{l+s+s1}{\PYZsq{}}\PY{p}{)}
\PY{n+nb}{print}\PY{p}{(}\PY{n}{urls}\PY{p}{)}
\end{Verbatim}
\end{tcolorbox}

    \begin{Verbatim}[commandchars=\\\{\}]
['https://gwosc.org/eventapi/json/GWTC-1-confident/GW170817/v3/H-
H1\_GWOSC\_4KHZ\_R1-1187008867-32.hdf5', 'https://gwosc.org/eventapi/json/GWTC-1-
confident/GW170817/v3/H-H1\_GWOSC\_4KHZ\_R1-1187006835-4096.hdf5', 'https://gwosc.o
rg/eventapi/json/GWTC-1-confident/GW170817/v3/L-L1\_GWOSC\_4KHZ\_R1-1187008867-
32.hdf5', 'https://gwosc.org/eventapi/json/GWTC-1-confident/GW170817/v3/L-
L1\_GWOSC\_4KHZ\_R1-1187006835-4096.hdf5', 'https://gwosc.org/eventapi/json/GWTC-1-
confident/GW170817/v3/V-V1\_GWOSC\_4KHZ\_R1-1187008867-32.hdf5', 'https://gwosc.org
/eventapi/json/GWTC-1-confident/GW170817/v3/V-V1\_GWOSC\_4KHZ\_R1-1187006835-
4096.hdf5', 'https://gwosc.org/eventapi/json/GWTC-1-confident/GW170817/v3/G-
G1\_GWOSC\_4KHZ\_R1-1187008867-32.hdf5', 'https://gwosc.org/eventapi/json/GWTC-1-
confident/GW170817/v3/G-G1\_GWOSC\_4KHZ\_R1-1187006835-4096.hdf5']
    \end{Verbatim}

    \begin{tcolorbox}[breakable, size=fbox, boxrule=1pt, pad at break*=1mm,colback=cellbackground, colframe=cellborder]
\prompt{In}{incolor}{17}{\boxspacing}
\begin{Verbatim}[commandchars=\\\{\}]
\PY{n}{urls} \PY{o}{=} \PY{n}{get\PYZus{}event\PYZus{}urls}\PY{p}{(}\PY{l+s+s1}{\PYZsq{}}\PY{l+s+s1}{GW170817}\PY{l+s+s1}{\PYZsq{}}\PY{p}{,} \PY{n}{duration}\PY{o}{=}\PY{l+m+mi}{4096}\PY{p}{,} \PY{n}{detector}\PY{o}{=}\PY{l+s+s1}{\PYZsq{}}\PY{l+s+s1}{V1}\PY{l+s+s1}{\PYZsq{}}\PY{p}{)}
\PY{n+nb}{print}\PY{p}{(}\PY{n}{urls}\PY{p}{)}
\end{Verbatim}
\end{tcolorbox}

    \begin{Verbatim}[commandchars=\\\{\}]
['https://gwosc.org/eventapi/json/GWTC-1-confident/GW170817/v3/V-
V1\_GWOSC\_4KHZ\_R1-1187006835-4096.hdf5']
    \end{Verbatim}

    Al realizar el filtro de la duración de los 4096 segundos y el detector
V1 para el evento GW170817 tenemos que el url del archivo que contiene
los datos es:

{[}`https://gwosc.org/eventapi/json/GWTC-1-confident/GW170817/v3/V-V1\_GWOSC\_4KHZ\_R1-1187006835-4096.hdf5'{]}

    \hypertarget{respuestas-a-preguntas-planteadas}{%
\section{Respuestas a preguntas
planteadas}\label{respuestas-a-preguntas-planteadas}}

    


    % Add a bibliography block to the postdoc
    
    
    
\end{document}
